\section{Introduction}
\label{sec:introduction}

Cross-cultural music recommendation becomes most difficult exactly where users value it most: when they want music that is unfamiliar in style but still meaningful in mood, function, or listening context. In this setting, a recommender must do more than retrieve acoustically similar tracks; it must bridge across repertoires without collapsing to the safest and most culturally dominant candidates. This matters for both music discovery and responsible MIR, because systems that ignore cultural asymmetry can systematically under-expose minority repertoires even when their upstream audio embeddings are strong \cite{holzapfel2018ethical,huang2023beyonddatasets,porcaro2021diversity}.

Recent music foundation models have substantially improved audio representation learning \cite{li2023mert,kanatas2025culturemert}, but current cross-cultural practice still leaves three gaps. First, embedding-only retrieval implicitly assumes that acoustic proximity is a sufficient proxy for cross-cultural relevance, even though unfamiliar style often appears as semantic distance in world-music settings \cite{kanatas2025culturemert,papaioannou2025universal}. Second, most current pipelines treat the backbone as the main locus of improvement, while the downstream stack for factorization, calibration, reranking, and feedback remains under-specified \cite{li2023mert,kanatas2025culturemert}. Third, standard recommender evaluation still centers relevance-style ranking metrics, even though cross-cultural deployment requires explicit control over cultural calibration, minority exposure at k, and data-source confound risk \cite{porcaro2021diversity,lee2025globalmood,huang2023beyonddatasets}. These gaps make it difficult to tell whether a system truly supports cross-cultural discovery or simply replays the geometry of a culturally imbalanced embedding space.

We therefore treat cross-cultural recommendation as a structured downstream problem rather than a pure backbone problem. Instead of proposing a new foundation model, we build a reusable stack on top of heterogeneous embedding backbones and make three contributions:
\begin{itemize}
\item We present a backbone-agnostic framework for culturally calibrated music recommendation that couples a unified V4 data contract with factorized downstream representation learning, stage-3 curriculum training, and calibration-aware reranking.
\item We show that calibration-aware reranking provides a controllable operating point between serendipity, cultural calibration, and minority exposure at k, and that this trade-off is observable on both CultureMERT and Gemini backbones.
\item We operationalize a participatory active learning pipeline that converts uncertainty-ranked candidate pairs into executable annotation packets and warm-start retraining hooks, making the system PAL-ready rather than purely offline.
\end{itemize}

The rest of the paper first reviews cross-cultural MIR, disentanglement, reranking, and human-feedback literature, then formalizes the V4 data contract and the DCAS framework, then describes the benchmark protocol on the V4 main dataset and the routeA_small sanity-check track, then presents main results together with module ablations and calibration sensitivity analyses, and finally discusses source confound, dataset coverage limits, and future work.
